\begin{claim}{AC.28}{Proved; mixed witnesses and exact prime/field minima}
For the pointed split curves $(C_c,\infty)$ of AC.25, the least
odd prime of a constant pointed L1 witness over all finite extensions
is eleven. A degree-eleven witness is defined over $\mathbf F_{4^{60}}$,
and sixty is the least extension degree for that prime.
The least extension degree permitting any constant odd-group witness
is sixteen, where a degree-$113$ witness exists.

For the following eight $\mathbf F_4$ models, write
$c\in\{\alpha,\alpha^2\}$ and $\delta\in\{0,\alpha\}$:
\[
 y^2+y=x+\frac{c}{x^2+x}+\delta,\qquad
 y^2+y=x+\frac{c}{x^2+x+\alpha}+\delta.
\]
Each original rational point has an eleven witness over
$\mathbf F_{4^{60}}$ and a $113$ witness over $\mathbf F_{4^{16}}$.
Whenever the cubic test AC.22 fails, eleven is the least odd prime.
In particular this holds for every rational point of the split
models and for infinity on each quadratic-pole model.
\end{claim}
\begin{proof}
We first specify and independently check the positive certificates.
Use the field
\[
 K=\mathbf F_2[t]/(t^{120}+t^4+t^3+t+1)=\mathbf F_{4^{60}},
\]
with binary polynomial coefficients encoded by integers, and set
\[
 \alpha=931750399618260646238930031330849980.
\]
The model is $Y^2+(x^2+x)Y=x^5+x^3+\alpha x^2+\alpha x$.
The following reduced pair defines the degree-eleven character line:
\[
 \begin{aligned}
 u(x)&=x^2+u_1x+u_0,\qquad v(x)=v_1x+v_0,\\
 u_0&=1286551050156223861353970523064590862,\\
 u_1&=855838071379080258553061447214623248,\\
 v_0&=372422475677710747247258963726964873,\\
 v_1&=880561393008999941135507307093448287.
 \end{aligned}
\]
The certificate verifies irreducibility of the field polynomial,
$\alpha^2+\alpha=1$, validity of the reduced pair, and exact order
eleven of its Jacobian class $a$. Since $q=2^{120}\equiv1\bmod11$,
the descent condition $\sigma_q(a)=qa$ holds over $K$.

All ten reduced multiples and Cartier coefficients are recorded in
\path{bielliptic_mixed_ell11_F4pow60.json}.
The independent verifier \path{verify_mixed_certificates.py} uses
pure Python carryless field arithmetic, with neither the searcher's
compiled field helper nor its Cartier-coefficient routine.
It recomputes every multiple and reconstructs doubling interpolation
from $(f+hv+v^2)/u$ and $h^{-1}\bmod u$.
All ten $q$-Frobenius scalars differ from one. Their product
determinant has the especially compact form
\[
          \overline{P_D(T)/P_C(T)}=1+T^4+T^{10}.
\]
Equivalently, the first scalar has reciprocal minimal polynomial
$1+T^2+T^5$ over $\mathbf F_2$, and its five conjugates each occur
twice. The determinant evaluates to one at $T=1$.
AC.27 therefore proves the odd relative class number over $K$,
with the cover pointed at infinity. This certificate determines
the factors modulo two; their full characteristic-zero coefficients
are not reconstructed here.

For the smaller field degree, the second certificate uses
\[
 E=\mathbf F_2[t]/(t^{224}+t^{21}+t^7+t+1)
\]
and its subfield $L=\mathbf F_{2^{32}}=\mathbf F_{4^{16}}$.
The complete reduced pair and all $112$ multiples are in
\path{bielliptic_mixed_ell113_F4pow16.json}.
They have exact order $113$ and satisfy $\sigma_{2^{32}}(a)=16a$,
where $16\equiv2^{32}\bmod113$.
Thus the constant pointed cover descends to $L$, although its
displayed character line uses the splitting field $E$.
The independent verifier checks all $112$ coefficients, applies
the actual inverse-$q$ conjugation in AC.27, and obtains an
$\mathbf F_2$ determinant with $31$ nonzero coefficients, whose
value at one is one. This proves the claimed witness over $L$.
Together with AC.26 it proves exact minimality of the field degree
sixteen.

AC.26 excludes three and five over every extension, and AC.30
excludes seven. The eleven witness therefore proves exact minimality
of the prime eleven. To minimize its field, the two elliptic factors
have modulo-eleven eigenvalue orders five and sixty respectively:
the first polynomial factors as $(X-3)(X-5)$, and the second has
both roots of exact order sixty. These facts are verified in
\path{bielliptic_obstructions.py}.
An eleven character using only one elliptic factor is excluded by
AC.24; using both requires an extension degree divisible by sixty.
Our witness attains it.

We also check transfer to the stated forms and points.
A quadratic-pole model becomes split after $x=z+s$, with
$s^2+s=\alpha$, followed by an Artin--Schreier constant change of
$y$. Both changes are available over $\mathbf F_{4^4}$ and preserve
infinity. The split quadratic twist is handled by the constant
change alone. This field embeds in both witness fields, of degrees
sixty and sixteen, so the transported witnesses retain exactly
those fields. Coefficientwise squaring transfers the certificates
from $c=\alpha$ to $c=\alpha^2$; no assertion that these two
geometric classes are $\mathbf F_4$-linearly isomorphic is needed.

The original Jacobian orders of these eight forms are $32,12,48,8$,
as given by their two point counts or the root's census. They are
prime to eleven and $113$. Hence any original point difference
$[Q-\infty]$ has order prime to the cover degree, also after transport
to the witness field. Every character of that prime degree kills
the difference. AC.6 proves that the transported cover is split
over every original $Q$, with the same relative polynomials.

For minimality, a hypothetical three, five or seven witness on a
form can be extended to the composite with $\mathbf F_{4^4}$.
The extra field degree is a power of two, which preserves the
numerical condition by AC.29. The isomorphism then transfers it
to the split model. Five and seven are unavailable at all original
points: their differences from infinity have orders dividing
$32,12,48$ or $8$, hence are divisible by those primes in every
relevant rational Jacobian, so their base-point twists are trivial.
For three we use the exact test AC.22. The split untwisted
Jacobians have two-power order, while the three points on each
split quadratic twist are hyperelliptic ramification points;
their cubic test fails. Infinity has $\kappa=0$ on every form.
This proves all minimal-prime assertions in the claim.
\end{proof}

\begin{claim}{AC.29}{Proved; safe extensions of characteristic-power degree}
Let $B/k$ be an abelian variety over a finite field of characteristic
$p$, and let $k'/k$ have degree $p^r$. Then
\[
 B(k)[p]=0\quad\Longleftrightarrow\quad B(k')[p]=0.
\]
In particular a relative-class-number $p$-unit for a fixed finite
separable cover is preserved by every constant extension of
$p$-power degree.
If two pointed curves become isomorphic over a $p$-power extension,
their existence of prime-degree witnesses over arbitrary extensions,
and the set of primes furnishing such witnesses, agree.
\end{claim}
\begin{proof}
On the finite-dimensional vector space
$V=B[p](\overline k)$ let $F$ be arithmetic Frobenius.
Rational $p$-torsion is its fixed space, and
\[
                     F^{p^r}-1=(F-1)^{p^r}.
\]
A power of an endomorphism of a finite-dimensional vector space
is invertible exactly when that endomorphism is invertible.
Thus the two fixed spaces are simultaneously zero.
Apply this to the connected new Jacobian factor to get the
relative-class-number assertion.
For pointed forms, the composite of any witness field with the
given $p$-power isomorphism field has $p$-power degree over the
witness field. The witness survives this extension and transports
across the pointed isomorphism. Apply the same argument in the
other direction.
\end{proof}

\begin{claim}{AC.30}{Proved; global seventh-prime exclusion for the bielliptic classes}
For $(C_c,\infty)$ of AC.25, no constant pointed degree-seven
L1 witness exists over any finite extension of $\mathbf F_4$.
The same holds for its coefficient conjugate and every pointed
form isomorphic to it at infinity over a power-of-two extension.
\end{claim}
\begin{proof}
Modulo seven the Frobenius polynomial from AC.25 is
\[
                  (X-2)^2(X^2-X+4).
\]
On the first primary part write $F=2I+N$, $N^2=0$.
On the second part every nonzero vector has period twenty-four
and $F^{16}=2I$. Consequently $F^{112}=2I$ on all of $J_C[7]$.
These polynomial identities and both possible Jordan types are
verified by \path{bielliptic_obstructions.py}.

The locally exact theta divisor is Frobenius-stable and symmetric:
if $B^2\simeq\omega_C$, then
$-[R-B]=[\iota R-B]$. A nonzero theta point $a$ of order seven
would therefore force all six nonzero multiples of $a$ into theta,
because $\langle2,-1\rangle=(\mathbf Z/7)^\times$.
This is impossible by the genus-two difference bound.
For a fixed nonzero Jacobian class, there are at most two ordered
representations as $R-S$: equality of two such representations
gives linearly equivalent effective degree-two divisors; if they
are distinct, both are the unique hyperelliptic pencil.
Thus a subset of size $r$ in the intersection of a theta translate
with a cyclic group of order seven satisfies
$r(r-1)\le2(7-1)$. It cannot have $r=6$.
Theta has no nonzero seven-torsion, and the checked locally exact
differential sequence makes every geometric cyclic seventh-degree
Prym ordinary, of two-rank six.

For a nonzero quotient functional $v$ on $J_C[7]$, its period $m$
and a scalar power of Frobenius are exhausted by
\[
 \begin{array}{c|c|c}
 m&d&vF^d/v\\\hline
 3&1&2\\
 21&7&2\\
 24&16&2\\
 168&112&2.
 \end{array}
\]
The first two cases are supported on the $(X-2)$ primary part,
without or with a nonzero Jordan contribution.
The last two include the quadratic primary part.
The certificate checks all $2400$ nonzero vectors for each Jordan
type, as well as the displayed formulas.

The corresponding geometric cover has a pointed model over
$\mathbf F_{4^d}$ with deck Frobenius equal to two.
Its new two-torsion is the sum of the two nontrivial irreducible
$\mathbf F_2[\mathbf Z/7]$-modules, each a one-dimensional
$\mathbf F_8$-space. On each, this Frobenius is semilinear for
squaring, and its cube is the norm to $\mathbf F_2^\times$,
hence the identity. New two-torsion is therefore rational over
$\mathbf F_{4^{3d}}$.
Here $3d=m$ or $2m$. The first field with constant deck group is
$\mathbf F_{4^m}$, and AC.29 shows that rational two-torsion cannot
first appear only in its quadratic extension.
Its relative class number is already even over that first field
and remains even over every extension.
The uniqueness of the pointed descent identifies every eligible
form with this base change. AC.29 also proves the final statement
for power-of-two pointed twists.
\end{proof}

